\ifdefined\iscombined
	\lecturetitle{Requirements}
\else
\documentclass{article}
\usepackage{listings}
\usepackage{amsthm}
\usepackage{comment}
\usepackage{amsmath}
\usepackage{amsfonts}
\usepackage{sectsty}
\usepackage{hyperref}
\usepackage{fancyhdr}
\usepackage[top=1in,bottom=1in,left=1in,right=1in]{geometry}
\usepackage{float}
\usepackage{graphicx}
\usepackage{tabularx}
\usepackage{mathtools}

\date{
	\vspace{-.75em}
	{\large \today}
}

\title{
	\vspace*{-1.5em}
	{\Huge Lecture 4} \\
	\vspace{-1em}
	\rule{\textwidth/4}{.01em} \\
	\vspace{-.25em}
	{\Large Requirements}
	\vspace{-.75em}
}

\author{
	{\large Matthew Farah}
}

\hypersetup{
	colorlinks=true,
	linkcolor=blue,
	filecolor=magenta,
	urlcolor=blue,
	citecolor=blue,
	anchorcolor=blue
}

\fancyhead{}
\fancyhead[L]{\today}
\fancyhead[R]{Lecture 4 - Requirements}

\sectionfont{\fontsize{12}{12}\bfseries}
\subsectionfont{\fontsize{10}{12}\normalfont}
\subsubsectionfont{\fontsize{10}{12}\normalfont}

\begin{document}

\maketitle
\pagestyle{fancy}

\fi

\begin{itemize}
	\item Design is best defined as the process of embedding some function into a form.
	\item The form a design takes should be:
	\begin{enumerate}
		\item Dependable.
		\item Sustainable.
		\item `Harmonious' with its environment
	\end{enumerate}
	\item Functional requirements describe the functionality of the design.
	\begin{itemize}
		\item Describes the design's reason for existence.
	\end{itemize}
	\item Non-functional requirements describe the form that design should take.
	\begin{itemize}
		\item Describes how the design should be manifested or instantiated
	\end{itemize}
	\item Non-functional requirements can put specific requirements on the design's speed, security, ease-of-use, and so on.
	\item The environment and functional requirements of a design place constraints on its implementation.
\end{itemize}

\ifdefined\iscombined
	\newpage
\else
	\end{document}
\fi
